\documentclass[14pt,a4paper,oneside]{extbook}
\usepackage[utf8]{vietnam}
\usepackage{amsmath,amssymb,mathrsfs,fancyhdr,enumerate,multirow,makecell}
\usepackage{tkz-euclide,tikz-3dplot,tikz,tkz-tab}
\usetikzlibrary{arrows.meta,calc,intersections,angles,quotes,backgrounds,shapes.geometric,patterns,through,shadows,positioning}
\usepackage{yhmath}
\usepackage[top=0.7cm, bottom=0.7cm, left=1.0cm, right=1.0cm, headheight=16pt, headsep=3mm, footskip=5mm]{geometry}
\usepackage[hidelinks,unicode]{hyperref}
\usepackage{qrcode}
\usepackage[many]{tcolorbox}
\usepackage{fontawesome}

\pagestyle{fancy}
\fancyhf{}
\fancyhead[L]{\footnotesize\bfseries\color{teal!85!black}\faBook~SỔ TAY KIẾN THỨC TOÁN 4}
\fancyhead[R]{\footnotesize\bfseries\color{blue!85!black}Lớp Toán Thầy Lĩnh -- loptoanthaylinh.vn}
\fancyfoot[C]{\footnotesize\bfseries\color{gray!80!black}-- Sổ tay kiến thức: Bài 07. Đo góc, đơn vị đo góc: Độ --}
\renewcommand{\headrulewidth}{0.4pt}

\begin{document}

% --- BANNER TIÊU ĐỀ BÀI HỌC ---
\noindent
\begin{tcolorbox}[
  enhanced,
  colback=teal!8!white,
  colframe=teal!80!black,
  arc=3mm, boxrule=1pt,
  left=3mm, right=3mm, top=2mm, bottom=2mm,
  halign=center
]
  {\Large\bfseries\color{teal!90!black}\faBookmark\ BÀI 07: ĐO GÓC, ĐƠN VỊ ĐO GÓC: ĐỘ}\\[1.5pt]
  {\footnotesize\itshape\color{black!75} Khái niệm góc, đỉnh, cạnh, thước đo góc và cách đọc số đo góc chính xác}
\end{tcolorbox}

\vspace{2.5mm}

% --- KHỐI 1 & 2: LÝ THUYẾT & KĨ NĂNG ---
\noindent
\begin{minipage}[t]{0.49\textwidth}
  \begin{tcolorbox}[
    enhanced,
    title={\footnotesize\bfseries\color{white}\faCompass\ 1. Khái niệm Góc \& Đơn vị đo},
    colback=blue!2!white,
    colframe=blue!75!black,
    arc=2.5mm, boxrule=0.8pt, top=1.8mm, bottom=1.8mm, left=2mm, right=2mm
  ]
    \footnotesize
    $\bullet$ \textbf{Cấu tạo góc:} Góc gồm có \textbf{đỉnh} và \textbf{hai cạnh} xuất phát từ đỉnh đó.\\
    \textit{Ví dụ:} Góc đỉnh $O$, cạnh $OA$ và $OB$.\\
    $\bullet$ \textbf{Đơn vị đo góc là Độ}, kí hiệu là $^\circ$.\\
    \textit{Ví dụ:} $60^\circ$ (đọc là: sáu mươi độ), $90^\circ$ (chín mươi độ).\\
    $\bullet$ Dụng cụ dùng để đo góc là \textbf{thước đo góc} (thước đo độ hình bán nguyệt).
  \end{tcolorbox}
\end{minipage}%
\hfill
\begin{minipage}[t]{0.49\textwidth}
  \begin{tcolorbox}[
    enhanced,
    title={\footnotesize\bfseries\color{white}\faSliders\ 2. Quy trình 3 bước đo góc},
    colback=teal!2!white,
    colframe=teal!75!black,
    arc=2.5mm, boxrule=0.8pt, top=1.8mm, bottom=1.8mm, left=2mm, right=2mm
  ]
    \footnotesize
    $\bullet$ \textbf{Bước 1:} Đặt \textbf{tâm của thước} trùng với \textbf{đỉnh} của góc.\\
    $\bullet$ \textbf{Bước 2:} Xoay thước sao cho một cạnh của góc đi qua \textbf{vạch $0^\circ$} của thước.\\
    $\bullet$ \textbf{Bước 3:} Cạnh còn lại của góc đi qua vạch nào trên thước thì đó là \textbf{số đo} của góc cần đo.
  \end{tcolorbox}
\end{minipage}

\vspace{2.5mm}

% --- KHỐI 3: SƠ ĐỒ TƯ DUY ---
\noindent
\begin{tcolorbox}[
  enhanced,
  title={\footnotesize\bfseries\color{white}\faSitemap\ Sơ đồ tư duy: Cấu tạo góc và thao tác sử dụng thước đo góc},
  colback=yellow!2!white,
  colframe=teal!85!black,
  arc=3mm, boxrule=0.9pt, top=1.8mm, bottom=1.8mm, left=1.5mm, right=1.5mm,
  halign=center
]
  \begin{tikzpicture}[
    scale=0.92,
    every node/.style={transform shape},
    root/.style={rectangle, rounded corners=5pt, fill=teal!85!black, text=white, font=\bfseries\footnotesize, inner sep=4.5pt, align=center},
    nodeStyle/.style={rectangle, rounded corners=4pt, fill=blue!10, draw=blue!75!black, text=blue!90!black, font=\bfseries\scriptsize, inner sep=3.5pt, align=center}
  ]
    \node[root] (core) at (0, 0.95) {\faCompass\ GÓC VÀ ĐƠN VỊ ĐO GÓC};

    \node[nodeStyle] (n1) at (-4.8, 0) {\textbf{CẤU TẠO GÓC}\\[1.5pt]Đỉnh $O$ và hai cạnh $OA, OB$};
    \node[nodeStyle] (n2) at (0, -0.1) {\textbf{ĐƠN VỊ ĐO}\\[1.5pt]Độ ($^\circ$), Thước đo góc};
    \node[nodeStyle] (n3) at (4.8, 0) {\textbf{CÁCH ĐỌC SỐ ĐO}\\[1.5pt]Nhìn đúng vòng vạch $0^\circ$};

    \draw[teal!85!black, thick, -{Stealth[length=2mm]}] (core.west) -| (n1.north);
    \draw[teal!85!black, thick, -{Stealth[length=2mm]}] (core.south) -- (n2.north);
    \draw[teal!85!black, thick, -{Stealth[length=2mm]}] (core.east) -| (n3.north);
  \end{tikzpicture}
\end{tcolorbox}

\vspace{2.5mm}

% --- KHỐI 4: MẸO & BẪY ---
\noindent
\begin{minipage}[t]{0.54\textwidth}
  \begin{tcolorbox}[
    enhanced,
    title={\footnotesize\bfseries\color{white}\faEye\ Mẹo chọn đúng vòng số trên thước đo độ},
    colback=green!3!white,
    colframe=green!60!black,
    arc=2.5mm, boxrule=0.8pt, top=1.8mm, bottom=1.8mm, left=2.2mm, right=2.2mm
  ]
    \footnotesize
    $\bullet$ Thước đo góc thường có \textbf{2 vòng số}: Vòng trong và vòng ngoài (từ $0^\circ$ đến $180^\circ$).\\
    $\bullet$ \textbf{Quy tắc vàng:} Cạnh đáy của góc trùng với vạch số 0 ở vòng nào (trong hay ngoài) thì \textbf{đọc số đo ở vòng đó}!\\
    $\bullet$ Góc bé hơn góc vuông chắc chắn $< 90^\circ$.
  \end{tcolorbox}
\end{minipage}%
\hfill
\begin{minipage}[t]{0.44\textwidth}
  \begin{tcolorbox}[
    enhanced,
    title={\footnotesize\bfseries\color{white}\faExclamationTriangle\ Mẹo tránh bẫy trắc nghiệm},
    colback=red!2!white,
    colframe=red!70!black,
    arc=2.5mm, boxrule=0.8pt, top=1.8mm, bottom=1.8mm, left=2.2mm, right=2.2mm
  ]
    \footnotesize
    $\bullet$ \textbf{Bẫy đọc nhầm số đo bù nhau:}\\
    Ví dụ góc nhọn $60^\circ$ lại đọc nhầm thành $120^\circ$ do nhìn sang vòng số đối diện!\\
    $\bullet$ \textbf{Bẫy đặt lệch tâm:}\\
    Tâm thước phải đặt \textbf{chính xác vào đỉnh góc}, nếu đặt lệch thì số đo bị sai hoàn toàn.
  \end{tcolorbox}
\end{minipage}

\end{document}
